\documentclass[11pt,a4paper]{article}

%====================================================
% PRÉAMBULE
%====================================================
\usepackage[utf8]{inputenc}
\usepackage[T1]{fontenc}
\usepackage[french]{babel}
\usepackage{lmodern}
\usepackage{geometry}
\geometry{margin=2cm}
\usepackage{amsmath,amssymb}
\usepackage{array,tabularx}
\usepackage{enumitem}
\usepackage{fancyhdr}
\usepackage{tikz}
\usepackage{tcolorbox}
\tcbuselibrary{skins,breakable}
\usepackage{multicol}

\pagestyle{fancy}
\fancyhf{}
\lhead{\textbf{Devoir surveillé -- Mathématiques 6e}}
\rhead{\textbf{Nombres, fractions, géométrie, données}}
\cfoot{\thepage}

\setlength{\parindent}{0pt}
\setlength{\parskip}{0.45em}

\definecolor{bleu}{RGB}{25,70,140}
\definecolor{grisclair}{RGB}{245,247,250}

\newtcolorbox{consignebox}{
  colframe=bleu,
  colback=grisclair,
  boxrule=0.8pt,
  arc=2mm,
  title=\textbf{Consignes générales}
}

\newcommand{\pts}[1]{\hfill \textbf{[#1 pt]}}
\newcommand{\ptss}[1]{\hfill \textbf{[#1 pts]}}

\begin{document}

%====================================================
% PAGE DE TITRE
%====================================================

\begin{center}
    {\Large \textbf{Devoir surveillé de mathématiques}}\\[0.2cm]
    {\large \textbf{Niveau Sixième -- Sujet ambitieux}}\\[0.2cm]
    {\large \textbf{Nombres, calculs, fractions, grandeurs, géométrie, données, probabilités et proportionnalité}}\\[0.4cm]
    \textbf{Durée conseillée : 1 h 45 à 2 h}\\[0.2cm]
    \textbf{Total : 80 points}
\end{center}

\begin{consignebox}
\begin{itemize}
    \item La calculatrice est autorisée seulement lorsque le professeur l'indique.
    \item Toutes les réponses doivent être justifiées par un calcul, une phrase ou une propriété du cours.
    \item Les unités doivent être précisées lorsqu'elles sont nécessaires.
    \item Les figures doivent être soignées et codées.
    \item Les résultats décimaux seront donnés exactement lorsque c'est possible, sinon arrondis au dixième si demandé.
    \item Le soin, la clarté des raisonnements et la qualité de la rédaction seront pris en compte.
\end{itemize}
\end{consignebox}

\vspace{0.3cm}

\begin{center}
\renewcommand{\arraystretch}{1.3}
\begin{tabularx}{0.95\textwidth}{|c|X|c|}
\hline
\textbf{Exercice} & \textbf{Thème principal} & \textbf{Points} \\
\hline
1 & Applications directes : décimaux, conversions, périmètre, aire & 10 \\
\hline
2 & Données : enquête exotique sur une serre martienne & 9 \\
\hline
3 & Probabilités : expérience à deux épreuves & 10 \\
\hline
4 & Géométrie exigeante : triangle, cercle, médiatrice, justification & 12 \\
\hline
5 & Fractions : calculs, comparaisons, justification des règles & 10 \\
\hline
6 & Problème transversal : sortie scolaire et proportionnalité & 11 \\
\hline
7 & Exercice original : machine à instructions & 8 \\
\hline
8 & Problème long : mini-parc d'aventures & 10 \\
\hline
\multicolumn{2}{|r|}{\textbf{Total}} & \textbf{80} \\
\hline
\end{tabularx}
\end{center}

\newpage

%====================================================
\section*{Exercice 1 -- Applications directes \hfill 10 points}
%====================================================

On rappelle que les réponses doivent être justifiées.

\begin{enumerate}[label=\textbf{\arabic*)}]

\item On considère le nombre :
\[
48\,205,307.
\]
\begin{enumerate}[label=\alph*)]
    \item Donner le chiffre des dizaines. \pts{1}
    \item Donner le chiffre des dixièmes. \pts{1}
    \item Donner le chiffre des millièmes. \pts{1}
    \item Écrire ce nombre sous la forme d'une somme faisant apparaître la valeur de chaque chiffre non nul. \ptss{2}
\end{enumerate}

\item Ranger les nombres suivants dans l'ordre croissant :
\[
7,09 \qquad 7,9 \qquad 7,099 \qquad 7,901.
\]
Justifier en ajoutant des zéros si nécessaire. \ptss{2}

\item Convertir :
\begin{enumerate}[label=\alph*)]
    \item $3,4$ m en cm. \pts{1}
    \item $250$ cL en L. \pts{1}
\end{enumerate}

\item Un rectangle mesure $8$ cm de longueur et $3,5$ cm de largeur.
Calculer son aire, en précisant l'unité. \pts{1}

\end{enumerate}

\hfill \textbf{Total exercice 1 : 10 points}

%====================================================
\section*{Exercice 2 -- Une enquête dans une serre martienne \hfill 9 points}
%====================================================

Dans une serre installée sur Mars, une classe observe la croissance de plusieurs plantes.
Les élèves relèvent la hauteur de $40$ plantes après trois semaines.

\[
\begin{array}{|c|c|}
\hline
\textbf{Hauteur des plantes} & \textbf{Nombre de plantes} \\
\hline
[0\ ;\ 5[ \text{ cm} & 4 \\
\hline
[5\ ;\ 10[ \text{ cm} & 9 \\
\hline
[10\ ;\ 15[ \text{ cm} & 15 \\
\hline
[15\ ;\ 20[ \text{ cm} & 8 \\
\hline
[20\ ;\ 25[ \text{ cm} & 4 \\
\hline
\end{array}
\]

\begin{enumerate}[label=\textbf{\arabic*)}]

\item Vérifier que le tableau contient bien $40$ plantes. \pts{1}

\item Combien de plantes mesurent au moins $10$ cm ? Justifier. \ptss{2}

\item Quelle fraction des plantes mesure entre $10$ cm et $15$ cm ? \pts{1}

\item Donner cette fraction sous forme décimale puis en pourcentage. \ptss{2}

\item Les scientifiques disent : « Plus de la moitié des plantes mesurent au moins $10$ cm. »
Ont-ils raison ? Justifier soigneusement. \ptss{2}

\item Pourquoi peut-on représenter ces données par un histogramme ? \pts{1}

\end{enumerate}

\hfill \textbf{Total exercice 2 : 9 points}

%====================================================
\section*{Exercice 3 -- Probabilités : les cristaux colorés \hfill 10 points}
%====================================================

Une boîte contient trois jetons numérotés :
\[
1,\quad 2,\quad 3.
\]
Une autre boîte contient deux cristaux :
\[
\text{un cristal rouge } R \quad \text{et un cristal bleu } B.
\]

On réalise l'expérience suivante :
\begin{itemize}
    \item on tire au hasard un jeton numéroté ;
    \item puis on tire au hasard un cristal.
\end{itemize}

Chaque jeton a la même chance d'être tiré, et chaque cristal a la même chance d'être tiré.

\begin{enumerate}[label=\textbf{\arabic*)}]

\item Compléter le tableau des issues possibles. \ptss{2}

\[
\renewcommand{\arraystretch}{1.4}
\begin{array}{|c|c|c|}
\hline
 & R & B \\
\hline
1 & (1;R) & \phantom{(1;B)} \\
\hline
2 & \phantom{(2;R)} & \phantom{(2;B)} \\
\hline
3 & \phantom{(3;R)} & \phantom{(3;B)} \\
\hline
\end{array}
\]

\item Combien y a-t-il d'issues possibles au total ? \pts{1}

\item Quelle est la probabilité d'obtenir le cristal rouge ? \ptss{2}

\item Quelle est la probabilité d'obtenir un numéro pair ? \ptss{2}

\item On considère l'événement :
\[
A=\text{« obtenir un numéro différent de }3\text{ »}.
\]
Calculer la probabilité de l'événement contraire de $A$, puis en déduire la probabilité de $A$. \ptss{2}

\item Lors d'une simulation, on répète cette expérience $60$ fois et on obtient $34$ fois le cristal rouge.
La fréquence observée du cristal rouge est-elle exactement égale à la probabilité théorique ? Est-ce surprenant ? Répondre par une phrase. \pts{1}

\end{enumerate}

\hfill \textbf{Total exercice 3 : 10 points}

%====================================================
\section*{Exercice 4 -- Géométrie : le trésor du triangle \hfill 12 points}
%====================================================

On considère un triangle $ABC$ tel que :
\[
AB=6\text{ cm}, \qquad AC=6\text{ cm}, \qquad \widehat{BAC}=40^\circ.
\]

\begin{enumerate}[label=\textbf{\arabic*)}]

\item Quelle est la nature du triangle $ABC$ ? Justifier avec une propriété du cours. \ptss{2}

\item Que peut-on dire des angles $\widehat{ABC}$ et $\widehat{ACB}$ ? Justifier. \ptss{2}

\item Calculer les mesures des angles $\widehat{ABC}$ et $\widehat{ACB}$. \ptss{3}

\item Construire le triangle $ABC$ en vraie grandeur. \ptss{2}

\item On veut placer un point $T$ qui soit à égale distance des points $B$ et $C$.
Sur quelle droite doit se trouver le point $T$ ? Donner le nom de cette droite. \pts{1}

\item Expliquer pourquoi tout point de cette droite est à égale distance de $B$ et de $C$. \ptss{2}

\end{enumerate}

\hfill \textbf{Total exercice 4 : 12 points}

%====================================================
\section*{Exercice 5 -- Fractions : ce qu'on a le droit de faire \hfill 10 points}
%====================================================

Cet exercice demande de calculer, mais aussi d'expliquer pourquoi certaines règles sont correctes ou fausses.

\begin{enumerate}[label=\textbf{\arabic*)}]

\item Calculer :
\[
\frac{3}{4}\text{ de } 28.
\]
Écrire le calcul complet. \ptss{2}

\item Compléter l'égalité :
\[
\frac{5}{6}=\frac{\cdots}{18}.
\]
Justifier par une phrase. \ptss{2}

\item Comparer :
\[
\frac{7}{8} \quad \text{et} \quad \frac{5}{8}.
\]
Justifier. \pts{1}

\item Comparer :
\[
\frac{3}{4} \quad \text{et} \quad \frac{5}{6}.
\]
On pourra écrire les deux fractions avec le même dénominateur. \ptss{2}

\item Un élève écrit :
\[
\frac{2}{5}+\frac{1}{5}=\frac{3}{10}.
\]
A-t-il raison ? Corriger son erreur en expliquant la règle. \ptss{2}

\item Un autre élève écrit :
\[
\frac{2}{3}=\frac{4}{6}.
\]
A-t-il raison ? Expliquer pourquoi. \pts{1}

\end{enumerate}

\hfill \textbf{Total exercice 5 : 10 points}

%====================================================
\section*{Exercice 6 -- Problème transversal : la sortie au planétarium \hfill 11 points}
%====================================================

Une classe de sixième organise une sortie au planétarium.

Il y a $24$ élèves et $2$ adultes accompagnateurs.

Le billet coûte :
\begin{itemize}
    \item $7,50$ euros par élève ;
    \item $10$ euros par adulte.
\end{itemize}

Le bus coûte $168$ euros pour tout le groupe.

Au planétarium, une animation dure $1$ h $35$ min.
La classe arrive à $13$ h $40$.

\begin{enumerate}[label=\textbf{\arabic*)}]

\item Calculer le prix total des billets pour les élèves. \pts{2}

\item Calculer le prix total des billets pour les adultes. \pts{1}

\item Calculer le coût total de la sortie, bus compris. \ptss{2}

\item On décide que chaque élève paiera la même somme, et que les adultes ne paient rien.
Calculer le prix à payer par élève. Donner une valeur exacte en euros. \ptss{2}

\item À quelle heure l'animation se termine-t-elle ? \ptss{2}

\item Le bus repart à $15$ h $30$.
La classe aura-t-elle au moins $10$ minutes pour aller aux toilettes et rejoindre le bus après l'animation ?
Justifier. \ptss{2}

\end{enumerate}

\hfill \textbf{Total exercice 6 : 11 points}

%====================================================
\section*{Exercice 7 -- Exercice original : la machine à instructions \hfill 8 points}
%====================================================

Une machine transforme un nombre de départ en suivant toujours les mêmes instructions :

\[
\boxed{
\text{Choisir un nombre}
\quad\longrightarrow\quad
\text{le multiplier par }4
\quad\longrightarrow\quad
\text{ajouter }6
\quad\longrightarrow\quad
\text{diviser par }2
}
\]

\begin{enumerate}[label=\textbf{\arabic*)}]

\item Quel résultat obtient-on si le nombre de départ est $5$ ? \ptss{2}

\item Quel résultat obtient-on si le nombre de départ est $12$ ? \ptss{2}

\item La machine affiche $19$.
Retrouver le nombre de départ en remontant les instructions à l'envers. \ptss{3}

\item Un élève affirme : « Cette machine donne toujours un nombre plus grand que le nombre de départ. »
A-t-il forcément raison ? Tester avec le nombre $1$ et conclure. \pts{1}

\end{enumerate}

\hfill \textbf{Total exercice 7 : 8 points}

%====================================================
\section*{Exercice 8 -- Problème long : le mini-parc d'aventures \hfill 10 points}
%====================================================

Un mini-parc d'aventures possède une zone rectangulaire de longueur $18$ m et de largeur $12$ m.

On veut installer :
\begin{itemize}
    \item une clôture tout autour de la zone, sauf une entrée de $2$ m ;
    \item un sol souple sur toute la zone ;
    \item des plots de jeu répartis dans la zone.
\end{itemize}

Le sol souple coûte $14$ euros par m$^2$.

Un plot occupe $3$ m$^2$.
Pour des raisons de sécurité, on ne peut utiliser que les $\dfrac{2}{3}$ de la surface totale pour installer les plots.

\begin{enumerate}[label=\textbf{\arabic*)}]

\item Calculer le périmètre complet de la zone rectangulaire. \pts{1}

\item Calculer la longueur de clôture nécessaire en tenant compte de l'entrée. \pts{1}

\item Calculer l'aire totale de la zone. \ptss{2}

\item Calculer le prix du sol souple. \ptss{2}

\item Calculer la surface maximale que l'on peut utiliser pour installer les plots. \ptss{2}

\item Combien de plots peut-on installer au maximum ? Justifier pourquoi il faut donner un nombre entier de plots. \ptss{2}

\end{enumerate}

\hfill \textbf{Total exercice 8 : 10 points}

\end{document}